\introduction % Структурный элемент: ВВЕДЕНИЕ

Современные организации придерживаются стратегии накопления всех доступных данных в
расчёте на их будущее использование в задачах, которые на момент сбора ещё не
сформулированы. Эта практика породила потребность в дешёвом, масштабируемом и при
этом аналитически пригодном хранилище. Классические хранилища данных (data
warehouse) обеспечивают транзакционность, управление схемой и быстрый SQL-доступ, но
плохо масштабируются по стоимости и не приспособлены к слабоструктурированным
данным. Озёра данных (data lake) на основе распределённых файловых систем и
объектных хранилищ дёшевы и масштабируемы, но не дают транзакционных гарантий и
управления метаданными. Архитектура lakehouse~\cite{armbrust2021lakehouse}
объединяет достоинства обоих подходов: поверх дешёвого объектного хранилища строится
слой открытого табличного формата, обеспечивающий ACID-транзакции, эволюцию схемы,
моментальные снимки и параллельный доступ. К числу таких систем относятся Apache
Iceberg~\cite{iceberg-spec}, Delta Lake~\cite{armbrust2020delta}, Apache
Hudi~\cite{hudi-docs} и Apache Paimon~\cite{paimon-docs}.

При проектировании lakehouse-хранилища одним из ключевых решений является выбор
файлового формата, в котором физически хранятся данные. От него зависят степень
сжатия (а значит, объём хранилища и стоимость), скорость записи, скорость
аналитических запросов и эффективность пропуска нерелевантных данных. Де-факто
стандартом колоночного хранения в экосистеме больших данных является Apache Parquet;
широко применяется также Apache ORC. Однако оба формата разрабатывались более
десяти лет назад, их форматы кодирования отражают возможности процессоров того
времени и не используют современные приёмы --- векторизованное декодирование с
явной SIMD-параллельностью, адаптивное сжатие чисел с плавающей точкой, кодирование
строк по схожести префиксов. Эти приёмы реализованы в новом колоночном формате
Vortex~\cite{vortex-repo, vortex-docs}, разрабатываемом с прицелом на нативное
(написанное на Rust) векторизованное выполнение и переносимость метаданных. Вопрос
о том, в каких сценариях преимущества
Vortex проявляются на практике и можно ли интегрировать его в существующую
lakehouse-систему, на момент начала работы оставался открытым.

Актуальность работы определяется, во-первых, практической потребностью в
обоснованном выборе формата хранения при построении витрин данных на Apache Paimon,
а во-вторых, тем, что Vortex --- активно развивающийся формат, поддержка которого в
Apache Paimon отсутствовала, и интеграция вместе с экспериментальной оценкой
представляет самостоятельный интерес как для сообщества Apache Paimon, так и для
организации, в инфраструктуре которой выполнялась работа.

Цель работы --- определить область применимости колоночного формата Vortex
в lakehouse-системе Apache Paimon, разработать модуль его интеграции и выполнить
сравнительный анализ форматов хранения (Parquet, ORC, Vortex) на представительных
нагрузках.

Для достижения цели поставлены следующие задачи:

\begin{enumerate}
    \item провести анализ подходов к хранению аналитических данных --- от хранилищ
        данных (Hive Metastore, Impala) к архитектуре lakehouse и открытым табличным
        форматам;
    \item провести обзор колоночных файловых форматов (Parquet, ORC, Vortex, Lance)
        и обосновать исключение строчного формата Avro из сравнения;
    \item изучить архитектуру Apache Paimon --- LSM-дерево, таблицы с первичным
        ключом, слияние при чтении, --- и алгоритмы кодирования, сжатия и пропуска
        данных, применяемые форматами;
    \item разработать модуль интеграции файлового формата Vortex в Apache Paimon;
    \item провести функциональное тестирование полученного артефакта локально в
        контейнерах;
    \item развернуть на кластере окружение нагрузочного тестирования средствами
        Ansible;
    \item провести серию экспериментов и собрать метрики производительности;
    \item сформулировать рекомендации по применению форматов (опционально).
\end{enumerate}

Объект исследования --- колоночные файловые форматы хранения данных в
составе lakehouse-системы Apache Paimon. Предмет исследования --- алгоритмы
кодирования, сжатия и пропуска данных, реализуемые этими форматами, и их влияние на
производительность чтения и записи таблиц.

Методы исследования: анализ исходного кода и документации Apache Paimon,
Apache Flink и формата Vortex; разработка программного модуля на языках Java и Rust;
функциональное тестирование в контейнеризованном окружении; нагрузочное тестирование
на кластере с использованием стандартизированного набора запросов TPC-DS и
специально разработанного потокового сценария; статистическая обработка результатов
измерений (усреднение по повторным запускам, оценка разброса).

Практическая значимость. Разработанный модуль интеграции расширяет перечень
поддерживаемых Apache Paimon форматов хранения. Полученные количественные результаты
и рекомендации применимы при проектировании витрин данных и настройке
lakehouse-хранилищ. Выявленное архитектурное ограничение интерфейса форматов Apache
Paimon (отсутствие проталкивания агрегатов) указывает направление дальнейшего
развития системы.

Отчёт состоит из введения, трёх разделов, заключения, списка использованных
источников и приложений. Первый раздел содержит аналитический обзор: эволюцию систем
хранения, открытые табличные форматы, колоночные файловые форматы, архитектуру
Apache Paimon, алгоритмы кодирования и пропуска данных, а также постановку задачи.
Второй раздел посвящён разработке модуля интеграции формата Vortex, функциональному
тестированию, развёртыванию окружения нагрузочного тестирования, методике
экспериментов и обнаруженным в ходе работы проблемам. Третий раздел содержит
результаты экспериментов, их анализ по типам запросов и рекомендации по применению
форматов.
